\label{sec:environment}

Agents are tasked with competing in a two-team hide-and-seek game in a physics-based environment. The \textit{hiders} are tasked with avoiding line of sight from the \textit{seekers}, and the seekers are tasked with keeping vision of the hiders. There are objects scattered throughout the environment that the agents can grab and also lock in place. There are also randomly generated immovable rooms and walls that the agents must learn to navigate. Before the game of hide-and-seek begins, the hiders are given a \textit{preparation phase} where the seekers are immobilized, giving the hiders a chance to run away or change their environment.

There are no explicit incentives for agents to interact with objects in the environment; the only supervision given is through the hide-and-seek objective. Agents are given a \textit{team based} reward; hiders are given a reward of 1 if all hiders are hidden and -1 if any hider is seen by a seeker. Seekers are given the opposite reward, -1 if all hiders are hidden and +1 otherwise.
To confine agent behavior to a reasonable space, agents are penalized with a reward of -10 if they go too far outside of the play area (outside an 18 meter square).
An episode lasts 240 timesteps, the first 40\% of which are the preparation phase where all agents are given zero reward.

We simulate the environment in the MUJOCO physics engine \citep{todorov2012mujoco}. The world is populated with 1 to 3 hiders, 1 to 3 seekers, 3 to 9 movable boxes of which at least 3 are elongated, 2 movable ramps, and randomly generated static walls and rooms.
We also experiment with a simpler, less randomized environment described in Appendix~\ref{appendix:quadrant}. Agents observe the position, velocity, and size (in the case of the randomly shaped boxes) of objects and other agents. If entities are not in line-of-sight of the agent or not in a 135 degree cone in front of the agent, then they are masked out in the policy. Agents also have 30 range sensors arrayed evenly around them, similar to a lidar.
In addition, each agent observes its own team and other agents' teams as well as how much time is left in the preparation phase.

Agents are simulated as spherical objects and have 3 action types that can be chosen simultaneously at each time step. They may \textit{move} by setting a discretized force along their $x$ and $y$ axis and torque around their $z$-axis.
They have a single binary action to \textit{grab} objects, which binds the agent to the closest object while the action is enabled. Agents may also \textit{lock} objects in place with a single binary action. Objects may be unlocked only by agents on the team of the agent who originally locked the object. Agents may only grab or lock objects that are in front of them and within a small radius.